\documentclass[conference]{IEEEtran}
\IEEEoverridecommandlockouts
\usepackage{cite}
\usepackage{balance}
\usepackage{amsmath,amssymb,amsfonts}
\usepackage{graphicx}
\usepackage{textcomp}
\usepackage{xcolor}
\usepackage{booktabs}
\usepackage{algorithm}
\usepackage{algorithmic}
\usepackage{tikz}
\usetikzlibrary{positioning,arrows.meta,fit,calc}
\usepackage[hidelinks]{hyperref}
\input{nums.tex}
\newtheorem{remark}{Remark}

\begin{document}

\title{One-Hot-Preserving QAOA for Flood Shelter\\
Assignment: Pairwise Capacity Encoding and Angle Transfer}

%<AUTHORS>
\author{
\IEEEauthorblockN{Nandhakumar T}
\IEEEauthorblockA{\textit{Department of Information Technology}\\
\textit{Sona College of Technology}\\
Salem, Tamil Nadu, India\\
nandhakumartm93@gmail.com}
\and
\IEEEauthorblockN{Balakarthikeyan B S}
\IEEEauthorblockA{\textit{Department of Artificial Intelligence and Data Science}\\
\textit{Sona College of Technology}\\
Salem, Tamil Nadu, India\\
balasan2626@gmail.com}
\and
\IEEEauthorblockN{Gowthaman J}
\IEEEauthorblockA{\textit{Department of Artificial Intelligence and Data Science}\\
\textit{Sona College of Technology}\\
Salem, Tamil Nadu, India\\
gowthamanjanadhran@gmail.com}
\and
\IEEEauthorblockN{Abhishek S}
\IEEEauthorblockA{\textit{Department of Artificial Intelligence and Data Science}\\
\textit{Sona College of Technology}\\
Salem, Tamil Nadu, India\\
abisheksengodan18@gmail.com}
}
%</AUTHORS>

\maketitle

\begin{abstract}
Shelter assignment after a flood disruption is a small assignment problem with capacity inequalities that is often proposed for quantum optimization. We study the Quantum Approximate Optimization Algorithm (QAOA) on 12-qubit instances from a flood digital twin by exact statevector simulation. Capacity is encoded with pairwise penalties instead of slack qubits; the encoding is approximate, and its minimiser was the true optimum in \QuboExact\% of instances. A ring-XY mixer that keeps each group assigned to one shelter placed a mean probability of \SeedXYMean{} on the optimal plan at depth 4 over \NSeeds{} optimizer seeds, against at most \SeedXMax{} for the transverse-field mixer at the same CNOT count per layer. One angle set per depth, chosen among the angles trained on \NPaired{} development instances, was applied unchanged to \HoN{} held-out seeds of the same disruption pattern. Chosen by mean optimal-plan probability, which needs the optima of the development instances, it reached \SelPopt{} against \SelTrained{} for per-instance training. Chosen without optima, by mean energy or CVaR, it reached \SelEn{} and \SelCv, but fell below per-instance training in leave-one-out tests. Transfer failed without a shared energy scale. On \CityN{} instances on the road networks of \CityNCities{} Indian cities, fixed angles left the mean probability unchanged but raised the chance that 32 verified samples contain the optimum. Exhaustive search solves every instance in about a millisecond; we claim no quantum advantage.
\end{abstract}

\begin{IEEEkeywords}
QAOA, XY mixer, parameter transfer, QUBO, evacuation planning
\end{IEEEkeywords}

\section{Introduction}
When a flood closes roads and fills shelters, an emergency officer must decide again which shelter each evacuation group should go to, trading off travel time, flood exposure and the seats left. As a binary program this is a capacity-constrained assignment problem, the kind that is routinely encoded as a quadratic unconstrained binary optimization (QUBO) problem \cite{Lucas2014,Glover2019} and offered as an application for near-term quantum optimization \cite{Cadelina2024,Abbas2024}.

Prior proposals rarely address three questions. First, does the circuit respect the constraints or only penalise violations? The original QAOA \cite{Farhi2014} explores all $2^n$ bitstrings, most of which are not plans, whereas the quantum alternating operator ansatz \cite{Hadfield2019} keeps the state inside a feasible subspace \cite{WangZ2020,Fuchs2022,Bartschi2020}. Mixer studies treat equality constraints, while shelter capacity is an inequality. Second, must the angles be trained again after every disruption, and how should a reusable angle set be chosen? Third, how do the results compare with classical methods and with what a sampler plus a feasibility check returns?

We address them on instances generated by the digital twin of a flood decision-support prototype. Our contributions are:
\begin{enumerate}
\item A shelter-assignment QUBO in which capacity is encoded by pairwise products instead of slack qubits, with a cost shift and scaling that place all instances on one energy range; we measure how often the encoding is exact and show that it admits infeasible zero-penalty plans.
\item A comparison of transverse-field, ring-XY, Grover and feasible-subspace mixers under mean-energy and conditional value-at-risk (CVaR) training over \NSeeds{} optimizer seeds, with Holm-corrected tests and CNOT counts per layer.
\item An angle-transfer study with three selection rules, two of which need no optima, an ablation in which transfer failed without the shared energy scale, and evidence that the performance of trained angle sets concentrates across instances.
\item The expected optimality gap of the best verified plan among $N$ samples, and a comparison with exhaustive search, a greedy rule and simulated annealing.
\end{enumerate}
All results come from classical simulation; we make no claim of quantum speed-up.

\section{Related Work}
\textit{Formulations.} Lucas gave Ising forms for many NP-hard problems \cite{Lucas2014}, and Glover \textit{et al.} describe how assignment and capacity constraints become quadratic penalties, with slack variables for inequalities \cite{Glover2019}. Flood shelter allocation has been solved by integer programming \cite{Sherali1991} and by quantum annealing \cite{Cadelina2024}.

\textit{Mixers.} QAOA alternates a cost unitary and a mixing unitary for $p$ layers \cite{Farhi2014,Blekos2024}. Hadfield \textit{et al.} generalised the mixer so that it preserves hard constraints \cite{Hadfield2019}. Wang \textit{et al.} analysed ring and complete XY mixers for one-hot constraints \cite{WangZ2020}, Fuchs \textit{et al.} gave general constructions of constraint-preserving mixers \cite{Fuchs2022} and mixers restricted to a subspace \cite{Fuchs2024}, and Grover mixers move the complexity into the preparation of the feasible superposition \cite{Bartschi2020}. Applied studies pair XY mixers with warm starts for vehicle routing \cite{doCarmo2025} and build coverage-preserving mixers for rostering \cite{Gupta2026}.

\textit{Training and transfer.} Angles can be interpolated from depth $p$ to $p{+}1$ \cite{Zhou2020}, CVaR objectives raise the probability of good samples \cite{Barkoutsos2020}, and warm starts use classical solutions \cite{Egger2021}. The objective value at fixed angles concentrates across typical instances \cite{Brandao2018}, optimal angles transfer between random graphs \cite{Galda2021}, and for weighted MaxCut a rescaling of the weights lets one parameter set serve many instances \cite{Shaydulin2023}. On hardware, QAOA performance falls with size and depth \cite{Harrigan2021}, and practical advantage has not been shown \cite{Abbas2024}.

\textit{Positioning.} We do not propose a new mixer; the ring-XY mixer is that of \cite{WangZ2020}. The transfer studies \cite{Brandao2018,Galda2021,Shaydulin2023} use unconstrained MaxCut with the standard mixer. We add a capacity inequality kept as a pairwise penalty next to a one-hot-preserving mixer, compare selection rules with and without optima, and evaluate on held-out instances and other road networks with a verifier in the loop.

\section{Problem Formulation}
\subsection{Instances from the Digital Twin}
The twin models a corridor with three population zones, nine junctions and three shelters joined by 22 roads. Every road carries a travel time, a flood-depth index and a congestion level drawn from a seed. Shelters have capacities of 120, 90 and 70 people with an initial occupancy of 15--45\%, and groups have 8--22 people. Each instance applies the disruption used in the twin's routing benchmark: road R5 closes and Shelter A receives arrivals equal to half its capacity. For group $g$ and shelter $s$ we take the best open path under the twin's flood- and congestion-weighted edge cost and define
\begin{equation}
c_{gs}=T_{gs}+5H_{gs},
\end{equation}
where $T_{gs}$ is travel time in minutes and $H_{gs}$ the mean flood index along the path. The pair $(g,s)$ is \emph{forbidden} if no open path exists, the shelter has failed, or the group alone exceeds the remaining seats $R_s$. With $G$ groups, binary variables $x_{gs}$ give $n=3G$ qubits. We write $m_g$ for the number of people in group $g$. A plan is feasible if every group has exactly one shelter, no forbidden pair is used and
\begin{equation}
\textstyle\sum_g m_g x_{gs}\le R_s\quad\text{for each shelter } s.
\end{equation}

\subsection{QUBO with Pairwise Capacity}
An inequality normally needs slack qubits \cite{Glover2019}. We use a pairwise encoding instead. Let $\mathcal{O}=\{(g,h,s): m_g+m_h>R_s\}$ be the pairs of groups that cannot share shelter $s$ and $\mathcal{F}$ the forbidden pairs. With shifted costs $c'_{gs}=c_{gs}-\min_{s'}c_{gs'}$ and penalty weight $\mu$,
\begin{multline}
E(\mathbf{x})=\sum_{g,s}c'_{gs}x_{gs}+\mu\Big[\sum_g\Big(1-\sum_s x_{gs}\Big)^2\\
+\sum_{(g,s)\in\mathcal{F}}x_{gs}+\sum_{(g,h,s)\in\mathcal{O}}x_{gs}x_{hs}\Big].
\label{eq:qubo}
\end{multline}

\begin{remark}
With $\mu=1+\sum_g\max_s c'_{gs}$, every bitstring that breaks a one-hot, forbidden-pair or pairwise-capacity term has higher energy than every bitstring that breaks none of them, since $c'_{gs}\ge0$ bounds the cost of a bitstring without violations by $\mu-1$ and each violation adds at least $\mu$. The bound no longer holds for the reduced weights $\lambda\mu$ with $\lambda<1$ used in the penalty sweep of Section~\ref{sec:mixres}.
\end{remark}

On one-hot states the shift subtracts the constant $\sum_g\min_s c_{gs}$, so the optimal plan is unchanged. We divide the Hamiltonian by $\Delta=\mu-1$, which maps the energies of all penalty-free plans into $[0,1]$ for every instance. The pairwise encoding misses overloads by three or more groups whose pairs all fit, so a penalty-free plan can still be infeasible, and every sampled plan is re-checked exactly before use.

\section{QAOA Variants and Training}
For depth $p$ and angles $\boldsymbol\theta=(\boldsymbol\gamma,\boldsymbol\beta)$, QAOA prepares
\begin{equation}
|\psi(\boldsymbol\theta)\rangle=\prod_{k=1}^{p}e^{-i\beta_k B}\,e^{-i\gamma_k \hat E/\Delta}\,|\psi_0\rangle,
\end{equation}
where $\hat E$ is diagonal with entries $E(\mathbf{x})$. Substituting $x_q=(1-Z_q)/2$ turns (\ref{eq:qubo}) into an Ising Hamiltonian with one $R_Z$ rotation per qubit and one $ZZ$ rotation, two CNOTs, per non-zero coupling.

\textit{Transverse-field (X) mixer.} $B=\sum_q X_q$ and $|\psi_0\rangle=|+\rangle^{\otimes n}$; constraints act only through the penalties. A cost layer has 12 one-hot couplings and on average \OvPairs{} capacity couplings, \CnotX{} CNOTs; the mixer needs none.

\textit{Ring-XY mixer.} For each group $g$ with qubits $a_g,b_g,c_g$,
\begin{equation}
B_{\mathrm{XY}}=\sum_{g}\ \sum_{(u,v)\in\{(a_g,b_g),(b_g,c_g),(c_g,a_g)\}}\tfrac12\big(X_uX_v+Y_uY_v\big).
\label{eq:xy}
\end{equation}
Each term swaps amplitude between $|01\rangle$ and $|10\rangle$, so a one-hot group stays one-hot \cite{WangZ2020}. The start state $|W_3\rangle^{\otimes G}$, with $|W_3\rangle=(|100\rangle+|010\rangle+|001\rangle)/\sqrt3$, is the uniform superposition of the $3^G$ one-hot plans \cite{Bartschi2019}. The one-hot penalty is then constant and dropped; the three ring terms are applied in a fixed order. Each XY rotation compiles to two CNOTs \cite{Vatan2004}, so a layer has \CnotXY{} CNOTs, the same count as the X variant. We prepare $|W_3\rangle$ with an $R_Y$ rotation, a controlled Hadamard (one CNOT) and two CNOTs, which we checked numerically; this adds \CnotW{} CNOTs once.

\textit{Grover mixer.} $B_{\mathrm{GM}}=|\psi_0\rangle\langle\psi_0|$ with the same start state, so $e^{-i\beta B_{\mathrm{GM}}}=I+(e^{-i\beta}-1)|\psi_0\rangle\langle\psi_0|$ \cite{Bartschi2020}. Each layer needs the inverse preparation, a phase on $|0\rangle^{\otimes n}$ controlled by all $n$ qubits and the preparation. With the textbook ancilla construction \cite{Nielsen2010}, the controlled phase uses $2(n-2)$ Toffoli gates of six CNOTs each, one controlled phase and $n-2=\NAnc$ ancilla qubits, \CnotMCP{} CNOTs in all, so a Grover layer has \CnotGM{} CNOTs.

\textit{Feasible-subspace mixer.} As a reference, a Grover mixer starts from the uniform superposition of the fully feasible plans \cite{Bartschi2020,Fuchs2024}. Preparing it requires the feasible set, so it is a bound, not a practical circuit.

\textit{Objectives and training.} We minimise either $\langle\hat E\rangle/\Delta$ or the CVaR with $\alpha=0.1$ \cite{Barkoutsos2020}. Sorting bitstrings by energy, $E_{(1)}\le E_{(2)}\le\dots$, with probabilities $\pi_{(k)}$ and cumulative mass $\Pi_k=\sum_{j\le k}\pi_{(j)}$,
\begin{equation}
\mathrm{CVaR}_\alpha=\frac1\alpha\sum_k E_{(k)}\,\min\!\big(\pi_{(k)},\ \max(0,\ \alpha-\Pi_{k-1})\big).
\label{eq:cvar}
\end{equation}
Every variant is trained from scratch with the same procedure: COBYLA \cite{Powell1994} with at most 250 evaluations per start, depth by depth. At each depth $p$ we start from two random points ($\gamma_k\sim U(0,\pi)$, $\beta_k\sim U(0,\pi/2)$) and, for $p>1$, from the optimum of depth $p-1$ interpolated to $p$ layers \cite{Zhou2020}, and keep the start with the lowest objective. Training never uses the optimal plan.

\subsection{Angle Transfer}\label{sec:transfer}
Since the scaled Hamiltonians of all instances share the energy range $[0,1]$, we test whether trained angles carry over. Let $\Theta_p=\{\boldsymbol\theta_j^{(p)}:j\in\mathcal{D}\}$ be the angles trained on the development instances $\mathcal{D}$. We fix one set $\boldsymbol\theta^{\star}_p\in\Theta_p$ by one of three rules:
\begin{equation}
\boldsymbol\theta^{\star}_p=\arg\max_{\boldsymbol\theta\in\Theta_p}\ \tfrac{1}{|\mathcal{D}|}\textstyle\sum_{j\in\mathcal{D}}P_{\mathrm{opt}}(j;\boldsymbol\theta),
\label{eq:transfer}
\end{equation}
or the set with the lowest mean of $\langle\hat E\rangle/\Delta$, or with the lowest mean CVaR, over $\mathcal{D}$. Rule (\ref{eq:transfer}) needs the exact optimum of every development instance, which we obtain by exhaustive search; it is an oracle on the development set, although it uses nothing from the instances it is applied to. The other two rules use only quantities that a device can estimate. We choose among whole trained sets because their coordinate-wise median, which mixes different basins of the periodic angle space, gave $P_{\mathrm{opt}}=\LooMedianFour$ in a leave-one-out test.

\section{Experimental Setup}\label{sec:setup}
\textit{Instances.} The development set has \NInst{} instances from twin seeds $42+17i$, $i=0,\dots,29$, with the first four groups ($G=4$, $n=12$). All are feasible, with on average \MeanFeasAssign{} feasible plans among the 81 one-hot plans, and in \NBinding{} of them capacity changes the optimal plan (\emph{capacity-binding}). The held-out set has \HoN{} further seeds, $i=30,\dots,59$, of the same disruption pattern, used only for the final transfer test. We also built \CityAll{} instances (\CityN{} feasible, \CityBinding{} capacity-binding) on the OpenStreetMap road graphs of \CityNCities{} flood-prone Indian cities, with each city's mapped zones, shelters and road flood-exposure proxy. Demand is not mapped, so group sizes and seats follow the twin; the disruption closes the most flood-exposed road on any zone--shelter path and halves the first shelter's seats.

\textit{Simulation.} The simulator stores all $2^n$ amplitudes in double precision; all random choices use fixed seeds.

\textit{Metrics.} $P_{\mathrm{opt}}$ is the probability of measuring an optimal plan and $P_{\mathrm{feas}}$ that of a feasible plan. A sampler draws $N$ plans and returns the cheapest one that passes an exact feasibility check. It returns the optimum whenever the optimum is among the samples, which has probability success@$N=1-(1-P_{\mathrm{opt}})^N$, assuming independent shots. We also report the expected relative gap $(c-c^\star)/c^\star$ of the returned plan, given that a feasible plan was drawn, for $N\in\{8,32,128\}$, computed exactly from the output distribution. One evaluation is one estimate of the training objective, which on hardware is a batch of circuit executions.

\textit{Statistics.} Paired comparisons use the two-sided Wilcoxon signed-rank test on per-instance values, also for success@$N$. We apply the Holm correction within each family of tests: the four depths of one mixer comparison, the \SelFamily{} fixed-versus-trained comparisons of Table~\ref{tab:select}, and the four city comparisons. Effect sizes are median paired differences. Intervals for variants run with \NSeeds{} optimizer seeds use a two-level bootstrap that resamples instances and then one seed per instance. Because the city instances cluster by city, we also apply a sign test to the six per-city mean differences.

\textit{Baselines.} Uniform sampling over the 81 one-hot plans ($P_{\mathrm{opt}}=\UniOneHot$) and over the feasible plans ($P_{\mathrm{opt}}=\UniFeas$); exhaustive search; a greedy rule that places the largest group first in its cheapest shelter with seats left; and simulated annealing (SA) \cite{Kirkpatrick1983} with 20 restarts of 200 sweeps, either with bit flips on (\ref{eq:qubo}) or with moves that reassign one group, which keep the plan one-hot. For the X mixer we swept the penalty $\lambda\mu$ with $\lambda\in\{0.25,0.5,1,2\}$.

\section{Results}
\subsection{Mixers, Objectives and Penalties}\label{sec:mixres}
Table~\ref{tab:mixers} and Fig.~\ref{fig:popt} give the comparison. Over \NSeeds{} seeds the XY mixer reached $P_{\mathrm{opt}}=\SeedXYMean$~$\pm$~\SeedXYSd{} at $p=4$ (seed means \SeedXYMin--\SeedXYMax), against \SeedXMin--\SeedXMax{} for the X mixer. It was better on \StXYvXWinsMin--\StXYvXWinsMax{} of \NPaired{} instances at every depth (Holm-adjusted $p\le\StXYvXMaxP$), and the ranking held for every seed. Its feasibility stayed between \XYPfeasMin{} and \XYPfeasMax{} at every depth, whereas the X mixer began with $P_{\mathrm{feas}}=\QXMeanOnePfeas$ and spent its layers recovering feasibility.

\begin{table}[t]
\caption{Mixers and objectives at depth 4 (\NPaired{} development instances, 12 qubits)}
\label{tab:mixers}
\centering
\footnotesize
\setlength{\tabcolsep}{2.6pt}
\begin{tabular}{lcccccc}
\toprule
Variant & $P_{\mathrm{opt}}$ & 95\% CI & $P_{\mathrm{feas}}$ & s@\Shots\,\% & Evals & CNOT\\
\midrule
\input{tab_mixers.tex}
\bottomrule
\multicolumn{7}{l}{\scriptsize X, XY: mean over \NSeeds{} seeds, CI by two-level bootstrap. Evals: objective}\\
\multicolumn{7}{l}{\scriptsize evaluations per instance, depths 1--4. CNOT: per layer, all-to-all connectivity.}\\
\end{tabular}
\end{table}

\begin{figure}[t]
\centering
\includegraphics[width=\columnwidth]{fig_popt.pdf}
\caption{Probability of the optimal plan (left) and of any feasible plan (right) against depth, optimizer seed 0. The two panels have different $y$-axis ranges. The dotted line is uniform sampling of one-hot plans.}
\label{fig:popt}
\end{figure}

The Grover mixer reached \GMMeanFour{} at $p=4$, and its difference from XY was not significant at any depth (Holm-adjusted $p\ge\StXYvGMMinP$). The data lean towards XY: at $p=4$ it was ahead on \GmWinsFour{} of \NPaired{} instances, with a mean difference of \GmDiffFour{} (paired bootstrap interval \GmDiffLo--\GmDiffHi) and an unadjusted $p=\GmRawP$. The Grover mixer spread probability more evenly over instances, so its success@\Shots{} was higher, but its layer costs \CnotGM{} CNOTs and \NAnc{} ancillas against \CnotXY{} CNOTs for ring-XY.

CVaR training lowered the XY mixer's mean $P_{\mathrm{opt}}$ but raised success@\Shots{} from \QXYMeanFourSucc\% to \QXYCvarFourSucc\% (seed 0): the worst instance rose from \XYMin{} to \XYCvarMin. The worst instances were all capacity-binding.

We tested whether a smaller penalty helps the X mixer: lowering the weight to $\lambda=\PenBestLam$ raised its $P_{\mathrm{opt}}$ by \PenGain$\times$ to \PenBestPopt{} at a feasibility of only \PenBestFeas, still about half of the XY result. To bound the cost of keeping capacity as a penalty, we ran a mixer restricted to feasible plans: it reached $P_{\mathrm{opt}}=\FgFour$ at $p=4$, \FgOverUni{} times uniform sampling over the same feasible plans, and beat XY on \FgWinsFour{} of \NPaired{} instances ($p=\FgPFour$).

\subsection{Angle Transfer to Held-Out Seeds}
\begin{table*}[t]
\caption{Angles at $p=4$ on the \HoN{} held-out seeds: per-instance training against one fixed angle set chosen on the development set}
\label{tab:select}
\centering
\footnotesize
\setlength{\tabcolsep}{4pt}
\begin{tabular}{llccccccc}
\toprule
Variant & Angles & Needs optima & Held-out $P_{\mathrm{opt}}$ & Median paired diff. & Wins of \HoN & Holm $p$ & Leave-one-out $P_{\mathrm{opt}}$ (dev.) & s@\Shots\,\%\\
\midrule
\input{tab_select.tex}
\bottomrule
\multicolumn{9}{l}{\scriptsize Fixed sets vs.\ per-instance training of the same variant (Wilcoxon signed-rank, Holm over \SelFamily{} rows). Per-instance rows: development-set mean in the leave-one-out column.}\\
\end{tabular}
\end{table*}

Table~\ref{tab:select} gives the held-out test. For the XY mixer with the mean objective, the set chosen by (\ref{eq:transfer}) reached $P_{\mathrm{opt}}=\SelPopt$ against \SelTrained{} for per-instance training with \HoCumFour{} evaluations (better on \SelPoptWins{} of \HoN, median difference \SelPoptMed, Holm $p=\SelPoptP$). The oracle-free sets also beat per-instance training on these seeds: \SelEn{} for the energy rule ($p=\SelEnP$) and \SelCv{} for the CVaR rule ($p=\SelCvP$). The leave-one-out test on the development set tells a different story. There, the energy rule gave \SelEnLoo{} and the CVaR rule \SelCvLoo, both below the \SelDevTrained{} of per-instance training, while rule (\ref{eq:transfer}) gave \SelPoptLoo. The same held for CVaR angles chosen by CVaR (\SelCvarCv{} on held-out seeds, \SelCvarCvLoo{} in leave-one-out against \SelCvarDevTrained). An oracle-free choice therefore did not give a reliable gain. Rule (\ref{eq:transfer}) gained in both tests; the optima it needs come offline from exhaustive search or integer programming at this size. Transfer did not rescue the X mixer (\SelXPopt).

Part of the gap is a mismatch of objectives. Per-instance training minimises energy well: its mean $\langle\hat E\rangle/\Delta$ was \EnTrained{} against \EnFixed{} for the fixed set, lower on \EnLower{} of \HoN{} instances. Low energy did not imply a high $P_{\mathrm{opt}}$, since across the 30 trained sets the rank correlation between mean energy and mean $P_{\mathrm{opt}}$ on the development set was only \RhoEnPopt. Stronger optimizers stayed below: the best of five COBYLA seeds by training objective reached \SelBestFive{} with \SelBestFiveEv{} evaluations per instance, and L-BFGS-B \SelLbfgs. COBYLA started from the fixed angles did improve on them (\StrWarm{} with \StrWarmEv{} evaluations, $p=\StrWarmP$).

Why does a fixed set work? The performance of trained sets concentrates across instances: a set's mean $P_{\mathrm{opt}}$ on the development set predicted its mean on the held-out seeds with rank correlation \RhoDevHeld. Only \CandBeat{} of the \CandN{} trained sets beat per-instance training there (median set \CandMedian), so the choice matters. The gain came mostly from capacity-binding instances: on the \HoNBind{} binding seeds per-instance training reached \HoBindTrained{} and the fixed set \HoBindFixed, on the others \HoNonBindTrained{} and \HoNonBindFixed. We repeated training, selection and the held-out test without the scaling and with division by the standard deviation of the penalty-free plan energies. Without scaling the fixed angles reached only \NormNoneFixed, although per-instance training still reached \NormNoneDevTrained; with the other normalisation they reached \NormStdFixed, close to our \NormDepFixed. Transfer needs a common energy scale, not our particular one. The one-time cost was \OffTrain{} evaluations of training and \OffCross{} for selection, the cost of per-instance training on about \OffBreakEven{} disruptions.

\subsection{Sampling with a Feasibility Check}
Table~\ref{tab:gap} gives what the sampler of Section~\ref{sec:setup} returns. With 8 samples, the fixed XY angles came within \GapFxEight\% of the optimal cost on average, per-instance training within \GapTrEight\% and uniform one-hot sampling within \GapUniEight\%. With 32 samples every XY variant drew a feasible plan with probability \PfeasThirtyTwoMin\% to one decimal, and with 128 samples per-instance training was within \GapTrBig\%. Many plans cost close to the optimum, so success@$N$ is the stricter measure.

\begin{table}[t]
\caption{Best verified plan among $N$ samples (held-out seeds, $p=4$): mean gap to the optimal cost and chance of drawing the optimum, in \%}
\label{tab:gap}
\centering
\footnotesize
\setlength{\tabcolsep}{2.2pt}
\begin{tabular}{lcccccc}
\toprule
 & \multicolumn{2}{c}{$N=8$} & \multicolumn{2}{c}{$N=32$} & \multicolumn{2}{c}{$N=128$}\\
\cmidrule(lr){2-3}\cmidrule(lr){4-5}\cmidrule(lr){6-7}
Sampler & gap & opt. & gap & opt. & gap & opt.\\
\midrule
\input{tab_gap.tex}
\bottomrule
\end{tabular}
\end{table}

\subsection{Transfer to Real Road Networks}
Table~\ref{tab:cities} breaks the city test down by city, with the angles fixed on the twin. Pooled over the \CityN{} instances, the set chosen by (\ref{eq:transfer}) reached $P_{\mathrm{opt}}=\CityFixed$ against \CityTrained{} for per-instance training (Holm $p=\CityPHolm$) and the energy-rule set \CityEnFixed{} ($p=\CityEnPHolm$). The fixed set led in only \CityPos{} of \CityK{} cities (sign test $p=\CitySignP$): no city-level gain in $P_{\mathrm{opt}}$. Per-instance training left \CityTrainedLow{} instances with $P_{\mathrm{opt}}<0.01$, the fixed set \CityFixedLow, and success@\Shots{} rose from \CitySuccTrained\% to \CitySuccFixed\% (Holm $p=\CitySuccPHolm$) and to \CityEnSucc\% for the energy rule ($p=\CityEnSuccPHolm$), in all \CitySuccPos{} cities (sign test $p=\CitySuccSignP$).

\begin{table}[t]
\caption{OpenStreetMap city instances, XY mixer, $p=4$: mean $P_{\mathrm{opt}}$ and success@\Shots\ (\%)}
\label{tab:cities}
\centering
\footnotesize
\setlength{\tabcolsep}{2.4pt}
\begin{tabular}{lccccccc}
\toprule
 & & \multicolumn{3}{c}{$P_{\mathrm{opt}}$} & \multicolumn{3}{c}{s@\Shots\,\%}\\
\cmidrule(lr){3-5}\cmidrule(lr){6-8}
City & inst. & trained & fix-O & fix-E & trained & fix-O & fix-E\\
\midrule
\input{tab_cities.tex}
\bottomrule
\multicolumn{8}{l}{\scriptsize fix-O: set chosen by max.\ mean $P_{\mathrm{opt}}$ (\ref{eq:transfer}); fix-E: by min.\ mean energy.}\\
\end{tabular}
\end{table}

\subsection{Larger Instances and Gate Noise}
We applied the 12-qubit angles of rule (\ref{eq:transfer}) unchanged to \SzFiveN{} new twin instances with five groups (15 qubits) and \SzSixN{} with six groups (18 qubits). They kept $P_{\mathrm{opt}}$ at least \SzFiveFixedMin{} at 15 qubits and \SzSixFixedMin{} at 18 qubits, \SzFiveRatio{} and \SzSixRatio{} times the mean of uniform one-hot sampling, while per-instance training fell below $0.01$ on \SzFiveTrainedLow{} and \SzSixTrainedLow{} instances. In mean $P_{\mathrm{opt}}$, \SizeCompare{} With ten instances per size, the test shows that fixed angles remain usable at 15 and 18 qubits; it does not establish a gain.

We simulated gate noise by quantum trajectories (\TrajRuns{} runs per circuit) on the first \TrajInst{} development instances: every two-qubit gate, including four gates per $W_3$ preparation, fails with probability $\epsilon$ and then applies one of the 15 non-identity two-qubit Pauli operators. Each $ZZ$ and XY gate compiles to two CNOTs, so $\epsilon=0.01$ corresponds to a CNOT error near 0.5\%, of the order of the \SycIso\% (isolated) and \SycSim\% (simultaneous) two-qubit Pauli errors reported for a superconducting processor \cite{Arute2019}. At $\epsilon=0.01$ and $p=4$ the XY mixer kept $P_{\mathrm{opt}}=\TrajXYOne$, \TrajRatioOne$\times$ the X mixer's \TrajXOne. A global depolarising model, in which a circuit with $N_2$ two-qubit gates returns its ideal state with probability $F=(1-\epsilon)^{N_2}$,
\begin{equation}
P_{\mathrm{opt}}^{\epsilon}=F\,P_{\mathrm{opt}}+(1-F)\,|\mathcal{X}^\star|/2^{n},
\label{eq:noise}
\end{equation}
with $\mathcal{X}^\star$ the optimal bitstrings, gave a lower \TrajGlobXYOne, since many Pauli errors barely change the output.

\subsection{Encoding Exactness and Classical Baselines}
In all \ZeroPenInfeasAll{} development instances some infeasible one-hot plan had zero penalty, because three groups whose pairs fit overloaded a shelter together. Such plans usually cost more than the optimum, so the minimiser of (\ref{eq:qubo}) was the true optimum in \QuboExact\% of instances. Exhaustive search took \ExactMs~ms per instance. The greedy rule found a feasible plan in \GreedyFeas\% of instances and the optimum in \GreedyOpt\%. Bit-flip SA ended at the optimum in \SATunedOpt\% of instances and SA with one-hot moves in \SAohOpt\%, the limit set by the encoding; both visited the optimum every time, so keeping the cheapest verified plan found it in \SAohVer\%.

\section{Discussion and Limitations}
Three findings are relevant for hardware experiments. First, the one-hot constraint belongs in the mixer: at the same CNOT count, ring-XY beat the transverse-field mixer on almost every instance and seed, and a smaller penalty did not close the gap. Capacity remains open: the feasible-subspace mixer roughly doubled $P_{\mathrm{opt}}$ but needs the feasible set. Second, transfer needs a common energy scale and a good selection rule. With development-set optima, one fixed set per depth matched or exceeded per-instance training at 12 qubits after a one-time offline cost of \OffTotal{} evaluations; without optima the gain was not reliable, and part of it reflects that energy is a weak proxy for $P_{\mathrm{opt}}$. Third, the encoding admits infeasible penalty-free plans, so a verifier must check every sample.

The study has clear limits. All held-out twin instances share one disruption pattern; other closures, failed shelters and capacity cuts were not tested, and the city instances use real roads but sampled demand. The size test stops at 18 qubits with ten instances per size; simulating the circuits on the $3^G$ one-hot amplitudes would allow more. The noise model charges errors per logical gate rather than per CNOT, ignores single-qubit and correlated errors and omits the Grover mixer, and no circuit was compiled for hardware. Absolute $P_{\mathrm{opt}}$ values depend on the training budget, and the Grover comparison cannot show equivalence. Future work will test varied disruptions, slack-variable encodings and hardware runs.

\section{Conclusion}
On flood shelter-assignment instances, a ring-XY mixer gave QAOA a significantly higher probability of the optimal plan than the transverse-field mixer at the same CNOT count, while capacity, kept as an approximate pairwise penalty, requires a verifier. Angles fixed by mean optimal-plan probability on development instances reached $P_{\mathrm{opt}}=\SelPopt$ on held-out seeds of the same disruption pattern, above per-instance training; rules without optima were not reliable, and transfer failed without a shared energy scale. Exhaustive search solves these instances in about a millisecond, so the results guide choices for hardware experiments and show no present-day advantage.

\section*{Acknowledgment}
We thank the OpenStreetMap contributors and Open-Meteo for open data used by the twin.
In line with the IEEE policy on AI-generated content, the authors disclose that Claude (Anthropic) was used to help edit and restructure the text of all sections and to draft the simulation, experiment and plotting scripts. All experiments are reproducible with fixed seeds. The authors checked every number, claim and reference and take full responsibility for the content.

\balance
\begin{thebibliography}{00}
\bibitem{Lucas2014} A. Lucas, ``Ising formulations of many NP problems,'' \textit{Front. Phys.}, vol. 2, Art. no. 5, 2014, doi: 10.3389/fphy.2014.00005.
\bibitem{Glover2019} F. Glover, G. Kochenberger, and Y. Du, ``Quantum bridge analytics I: A tutorial on formulating and using QUBO models,'' \textit{4OR}, vol. 17, no. 4, pp. 335--371, 2019, doi: 10.1007/s10288-019-00424-y.
\bibitem{Cadelina2024} A. R. Cadeli\~{n}a \textit{et al.}, ``D-Wave implementation of quantum annealing for optimal resource allocation in disaster response operation of Marikina City,'' \textit{ECTI Trans. Comput. Inf. Technol.}, vol. 18, no. 1, pp. 24--33, 2024, doi: 10.37936/ecti-cit.2024181.253577.
\bibitem{Abbas2024} A. Abbas \textit{et al.}, ``Challenges and opportunities in quantum optimization,'' \textit{Nat. Rev. Phys.}, vol. 6, pp. 718--735, 2024, doi: 10.1038/s42254-024-00770-9.
\bibitem{Farhi2014} E. Farhi, J. Goldstone, and S. Gutmann, ``A quantum approximate optimization algorithm,'' arXiv:1411.4028, 2014.
\bibitem{Hadfield2019} S. Hadfield, Z. Wang, B. O'Gorman, E. G. Rieffel, D. Venturelli, and R. Biswas, ``From the quantum approximate optimization algorithm to a quantum alternating operator ansatz,'' \textit{Algorithms}, vol. 12, no. 2, Art. no. 34, 2019, doi: 10.3390/a12020034.
\bibitem{WangZ2020} Z. Wang, N. C. Rubin, J. M. Dominy, and E. G. Rieffel, ``XY mixers: Analytical and numerical results for the quantum alternating operator ansatz,'' \textit{Phys. Rev. A}, vol. 101, no. 1, Art. no. 012320, 2020, doi: 10.1103/PhysRevA.101.012320.
\bibitem{Fuchs2022} F. G. Fuchs, K. O. Lye, H. M. Nilsen, A. J. Stasik, and G. Sartor, ``Constraint preserving mixers for the quantum approximate optimization algorithm,'' \textit{Algorithms}, vol. 15, no. 6, Art. no. 202, 2022, doi: 10.3390/a15060202.
\bibitem{Bartschi2020} A. B\"{a}rtschi and S. Eidenbenz, ``Grover mixers for QAOA: Shifting complexity from mixer design to state preparation,'' in \textit{Proc. IEEE Int. Conf. Quantum Comput. Eng. (QCE)}, 2020, pp. 72--82, doi: 10.1109/QCE49297.2020.00020.
\bibitem{Sherali1991} H. D. Sherali, T. B. Carter, and A. G. Hobeika, ``A location-allocation model and algorithm for evacuation planning under hurricane/flood conditions,'' \textit{Transp. Res. B, Methodol.}, vol. 25, no. 6, pp. 439--452, 1991, doi: 10.1016/0191-2615(91)90037-J.
\bibitem{Blekos2024} K. Blekos \textit{et al.}, ``A review on quantum approximate optimization algorithm and its variants,'' \textit{Phys. Rep.}, vol. 1068, pp. 1--66, 2024, doi: 10.1016/j.physrep.2024.03.002.
\bibitem{Fuchs2024} F. G. Fuchs and R. Pariente Bassa, ``LX-mixers for QAOA: Optimal mixers restricted to subspaces and the stabilizer formalism,'' \textit{Quantum}, vol. 8, Art. no. 1535, 2024, doi: 10.22331/q-2024-11-25-1535.
\bibitem{doCarmo2025} R. S. do Carmo, M. C. S. Santana, F. F. Fanchini, V. H. C. de Albuquerque, and J. P. Papa, ``Warm-starting QAOA with XY mixers: A novel approach for quantum-enhanced vehicle routing optimization,'' arXiv:2504.19934, 2025.
\bibitem{Gupta2026} A. Gupta and S. R. Hassan, ``Constraint-preserving QAOA for personnel rostering: Coverage-preserving and guarded-XY mixer constructions,'' arXiv:2607.09145, 2026.
\bibitem{Zhou2020} L. Zhou, S.-T. Wang, S. Choi, H. Pichler, and M. D. Lukin, ``Quantum approximate optimization algorithm: Performance, mechanism, and implementation on near-term devices,'' \textit{Phys. Rev. X}, vol. 10, no. 2, Art. no. 021067, 2020, doi: 10.1103/PhysRevX.10.021067.
\bibitem{Barkoutsos2020} P. K. Barkoutsos, G. Nannicini, A. Robert, I. Tavernelli, and S. Woerner, ``Improving variational quantum optimization using CVaR,'' \textit{Quantum}, vol. 4, Art. no. 256, 2020, doi: 10.22331/q-2020-04-20-256.
\bibitem{Egger2021} D. J. Egger, J. Mare\v{c}ek, and S. Woerner, ``Warm-starting quantum optimization,'' \textit{Quantum}, vol. 5, Art. no. 479, 2021, doi: 10.22331/q-2021-06-17-479.
\bibitem{Brandao2018} F. G. S. L. Brand\~{a}o, M. Broughton, E. Farhi, S. Gutmann, and H. Neven, ``For fixed control parameters the quantum approximate optimization algorithm's objective function value concentrates for typical instances,'' arXiv:1812.04170, 2018.
\bibitem{Galda2021} A. Galda, X. Liu, D. Lykov, Y. Alexeev, and I. Safro, ``Transferability of optimal QAOA parameters between random graphs,'' in \textit{Proc. IEEE Int. Conf. Quantum Comput. Eng. (QCE)}, 2021, pp. 171--180, doi: 10.1109/QCE52317.2021.00034.
\bibitem{Shaydulin2023} R. Shaydulin, P. C. Lotshaw, J. Larson, J. Ostrowski, and T. S. Humble, ``Parameter transfer for quantum approximate optimization of weighted MaxCut,'' \textit{ACM Trans. Quantum Comput.}, vol. 4, no. 3, Art. no. 19, 2023, doi: 10.1145/3584706.
\bibitem{Harrigan2021} M. P. Harrigan \textit{et al.}, ``Quantum approximate optimization of non-planar graph problems on a planar superconducting processor,'' \textit{Nat. Phys.}, vol. 17, pp. 332--336, 2021, doi: 10.1038/s41567-020-01105-y.
\bibitem{Bartschi2019} A. B\"{a}rtschi and S. Eidenbenz, ``Deterministic preparation of Dicke states,'' in \textit{Fundamentals of Computation Theory} (LNCS, vol. 11651). Cham, Switzerland: Springer, 2019, pp. 126--139, doi: 10.1007/978-3-030-25027-0\_9.
\bibitem{Vatan2004} F. Vatan and C. Williams, ``Optimal quantum circuits for general two-qubit gates,'' \textit{Phys. Rev. A}, vol. 69, no. 3, Art. no. 032315, 2004, doi: 10.1103/PhysRevA.69.032315.
\bibitem{Nielsen2010} M. A. Nielsen and I. L. Chuang, \textit{Quantum Computation and Quantum Information}, 10th anniversary ed. Cambridge, U.K.: Cambridge Univ. Press, 2010, doi: 10.1017/CBO9780511976667.
\bibitem{Powell1994} M. J. D. Powell, ``A direct search optimization method that models the objective and constraint functions by linear interpolation,'' in \textit{Advances in Optimization and Numerical Analysis}, S. Gomez and J.-P. Hennart, Eds. Dordrecht, The Netherlands: Kluwer, 1994, pp. 51--67, doi: 10.1007/978-94-015-8330-5\_4.
\bibitem{Kirkpatrick1983} S. Kirkpatrick, C. D. Gelatt, and M. P. Vecchi, ``Optimization by simulated annealing,'' \textit{Science}, vol. 220, no. 4598, pp. 671--680, 1983, doi: 10.1126/science.220.4598.671.
\bibitem{Arute2019} F. Arute \textit{et al.}, ``Quantum supremacy using a programmable superconducting processor,'' \textit{Nature}, vol. 574, no. 7779, pp. 505--510, 2019, doi: 10.1038/s41586-019-1666-5.
\end{thebibliography}

\end{document}
